\subsection{Restricciones}

\subsubsection{Regla: prohibido capturar o lanzar \texttt{Exception} genérica}

\textbf{Regla:} Prohibido lanzar directamente \texttt{System.Exception} u otro tipo genérico de excepción base, o capturarlo en un bloque \texttt{catch}; siempre se usa el tipo más específico. Esta restricción no impide heredar de \texttt{Exception} ni recibirla como parámetro \texttt{innerException} en el constructor convencional de una excepción personalizada.

\textbf{Descripción:} Regla de análisis \texttt{CA1031} de .NET: las excepciones genéricas ocultan problemas en tiempo de ejecución y dificultan la depuración (Microsoft, 2016).

\textbf{Código correcto:}
\begin{lstlisting}[style=csharpstyle]
public void EstablishConnection()
{
    try
    {
        _database.Connect();
    }
    catch (TimeoutException ex)
    {
        _logger.Error("Database connection timed out.", ex);
        throw;
    }
}
\end{lstlisting}

\textbf{Código incorrecto:}
\begin{lstlisting}[style=csharpstyle]
public void EstablishConnection()
{
    try
    {
        _database.Connect();
    }
    catch (Exception ex)
    {
        _logger.Error("Database connection failed.", ex);
        throw;
    }
}
\end{lstlisting}

\subsubsection{Regla: preferir excepciones predefinidas}

\textbf{Regla:} Se usa una excepción predefinida cuando expresa correctamente el error. Para argumentos nulos se usa \texttt{ArgumentNullException}; para argumentos inválidos, \texttt{ArgumentException} o una derivada; y para una operación incompatible con el estado actual, \texttt{InvalidOperationException}. Solo se crea una excepción personalizada cuando no existe una equivalente y el llamador necesita tratar esa condición por separado (Microsoft, 2025a).

\textbf{Código correcto:}
\begin{lstlisting}[style=csharpstyle]
public void StartTurn(Player player)
{
    ArgumentNullException.ThrowIfNull(player);

    if (!player.IsActive)
    {
        throw new InvalidOperationException("The player is not active.");
    }

    BeginTurn(player);
}
\end{lstlisting}

\textbf{Código incorrecto:}
\begin{lstlisting}[style=csharpstyle]
public void StartTurn(Player player)
{
    ArgumentNullException.ThrowIfNull(player);

    if (!player.IsActive)
    {
        throw new GameStateException("The player cannot start a turn.");
    }

    BeginTurn(player);
}
\end{lstlisting}
